\chapter*{Abstract}
\addcontentsline{toc}{chapter}{Abstract}
\label{ch:Abstract}

Traffic signal control at the network level constitutes a canonical multi-agent decision problem in which inter-agent coordination can substantially reduce urban congestion. Graph-based multi-agent reinforcement learning (MARL) methods, which treat signalized intersections as nodes in a communication graph, have shown promise in synthetic environments but remain underexplored on real-world German road networks characterized by irregular topologies and heterogeneous phase structures.

This thesis proposes FGS (FRAP-Graph-SAC), a MARL framework integrating three components: a phase-aware local encoder derived from FRAP, GATv2 graph-attention communication between neighboring intersections, and discrete-action Soft Actor-Critic (SAC) optimization. A second-generation variant, FGSv2, replaces the single node embedding with per-phase action tokens, enabling richer action-conditioned reasoning in both the actor and the centralized critic. The framework is implemented on top of SUMO, SUMO-RL, and RLlib, and evaluated under the centralized-training-with-decentralized-execution (CTDE) paradigm using the differential waiting time as the reward signal.

The framework is evaluated on all six German RESCO benchmark scenarios spanning three scale tiers — Single-intersection (cologne1, ingolstadt1), Corridor (cologne3, ingolstadt7), and Regional (cologne8, ingolstadt21) — against seven baselines including fixed-time control, max-pressure, IDQN, IPPO, independent SAC, FRAP, and CoLight. Performance is reported on mean waiting time, mean trip duration, mean delay, and throughput.

Results show that GATv2 communication yields the greatest benefit in Regional-tier scenarios, where heterogeneous network topology provides meaningful signal for attention weighting; gains are negligible in Single-intersection and Corridor settings. The FRAP encoder does not consistently outperform the MLP baseline owing to a structural mismatch between its phase-competition assumptions and German intersection geometries. SAC achieves competitive performance on medium-scale networks but encounters critic-scalability constraints at the largest evaluated scenario (ingolstadt21, $N=21$), where PPO-based variants are more reliable. The best-performing FGS configuration, \mbox{FGS(MLP+GATv2+SAC)}, attains a mean waiting time of 4.9\,s on the Regional cologne8 scenario, outperforming all baselines including CoLight (23.8\,s).

These findings demonstrate that graph-aware MARL with GATv2 communication is a viable and analytically informative framework for German-scenario traffic signal control, and identify the conditions under which phase-aware encoding and entropy-regularized optimization produce the greatest benefit.

\medskip
\noindent\textbf{Keywords:} multi-agent reinforcement learning, traffic signal control, graph attention network, soft actor-critic, FRAP encoder, RESCO benchmark, decentralized execution, German traffic scenarios

\chapter*{Kurzfassung}
\addcontentsline{toc}{chapter}{Kurzfassung}
\label{ch:Kurzfassung}

Die Verkehrssignalsteuerung auf Netzwerkebene stellt ein kanonisches Mehrfach-Agenten-Entscheidungsproblem dar, bei dem eine koordinierte Abstimmung zwischen Agenten den städtischen Verkehrsstau erheblich verringern kann. Graphbasierte Methoden des Multi-Agenten Reinforcement Learning (MARL), die Lichtsignalanlagen als Knoten in einem Kommunikationsgraphen modellieren, haben in synthetischen Umgebungen vielversprechende Ergebnisse gezeigt, wurden jedoch für reale deutsche Straßennetze mit unregelmäßigen Topologien und heterogenen Phasenstrukturen kaum untersucht.

Diese Arbeit schlägt FGS (FRAP-Graph-SAC) vor, ein MARL-Framework, das drei Komponenten integriert: einen phasenbewussten lokalen Kodierer auf Basis von FRAP, eine GATv2-Graphaufmerksamkeitskommunikation zwischen benachbarten Kreuzungen sowie eine diskrete Soft Actor-Critic (SAC)-Optimierung. Eine Weiterentwicklung, FGSv2, ersetzt die einzelne Knotenrepräsentation durch phasenspezifische Aktions-Token, die eine reichhaltigere aktionsbedingte Verarbeitung sowohl im Akteur als auch im zentralisierten Kritiker ermöglichen. Das Framework ist auf Basis von SUMO, SUMO-RL und RLlib implementiert und wird nach dem CTDE-Paradigma (Centralized Training with Decentralized Execution) mit der differenziellen Wartezeit als Belohnungssignal evaluiert.

Das Framework wird auf allen sechs deutschen RESCO-Benchmark-Szenarien in drei Skalierungsstufen evaluiert — Einzelkreuzung (cologne1, ingolstadt1), Korridor (cologne3, ingolstadt7) und Regional (cologne8, ingolstadt21) — und mit sieben Basisverfahren verglichen, darunter Festzeitsteuerung, Maximaldruck, IDQN, IPPO, unabhängiges SAC, FRAP und CoLight. Die Leistung wird anhand von mittlerer Wartezeit, mittlerer Fahrzeit, mittlerer Verzögerung und Durchsatz bewertet.

Die Ergebnisse zeigen, dass GATv2-Kommunikation den größten Nutzen in Szenarien der Regionalstufe erzielt, wo heterogene Netzwerktopologien aussagekräftige Signale für die Aufmerksamkeitsgewichtung liefern; in Einzelkreuzungs- und Korridorumgebungen sind die Verbesserungen vernachlässigbar. Der FRAP-Kodierer übertrifft die MLP-Baseline nicht konsistent, da eine strukturelle Diskrepanz zwischen den Phasenwettbewerbs-Annahmen und den deutschen Kreuzungsgeometrien besteht. SAC erzielt wettbewerbsfähige Leistungen in mittleren Netzwerken, stößt jedoch beim größten evaluierten Szenario (ingolstadt21, $N=21$) auf Skalierungsprobleme beim Kritiker, wo PPO-basierte Varianten zuverlässiger sind. Die beste FGS-Konfiguration, \mbox{FGS(MLP+GATv2+SAC)}, erreicht eine mittlere Wartezeit von 4.9\,s im regionalen Cologne-8-Szenario und übertrifft damit alle Basisverfahren einschließlich CoLight (23.8\,s).

Diese Erkenntnisse belegen, dass graphbewusstes MARL mit GATv2-Kommunikation ein praktikables und analytisch aufschlussreiches Framework für die deutsche Verkehrssignalsteuerung darstellt, und identifizieren die Bedingungen, unter denen phasenbewusstes Kodieren und entropieregularisierte Optimierung den größten Nutzen bringen.

\medskip
\noindent\textbf{Schlüsselwörter:} Multi-Agenten Reinforcement Learning, Verkehrssignalsteuerung, Graphaufmerksamkeitsnetzwerk, Soft Actor-Critic, FRAP-Kodierer, RESCO-Benchmark, dezentralisierte Ausführung, deutsche Verkehrsszenarien
